\lecture{05}{Выпуклые множества и выпуклые функции}

\sourcevideo
    {https://www.youtube.com/playlist?list=PLOzRYVm0a65fhKDS8cYIC7YCuVGJ4FOJx}
    {Week 1: Lecture 5: Convex Sets and Convex Functions}

Выпуклость обеспечивает глобальные гарантии в задачах оптимизации.
На выпуклой области всякий локальный минимум выпуклой функции является
глобальным. Для дифференцируемых функций выпуклость можно проверять по
градиенту, а для дважды дифференцируемых --- по матрице Гессе.

\section{Выпуклые комбинации и множества}

Пусть \(x,y\in\RR^n\).

\begin{definition}[Линейная комбинация]
Точка
\[
    z=\alpha x+\beta y,
    \qquad
    \alpha,\beta\in\RR,
\]
называется линейной комбинацией \(x\) и \(y\).
\end{definition}

\begin{definition}[Аффинная комбинация]
Линейная комбинация называется аффинной, если сумма коэффициентов
равна единице:
\[
    z=\alpha x+\beta y,
    \qquad
    \alpha+\beta=1.
\]
Для двух точек её можно записать как
\[
    z=\theta x+(1-\theta)y,
    \qquad
    \theta\in\RR.
\]
\end{definition}

\begin{definition}[Выпуклая комбинация]
Аффинная комбинация называется выпуклой, если её коэффициенты
неотрицательны:
\[
    z=\theta x+(1-\theta)y,
    \qquad
    \theta\in[0,1].
\]
\end{definition}

При изменении \(\theta\) от \(0\) до \(1\) точка \(z\) пробегает
отрезок, соединяющий \(x\) и \(y\).

Для точек \(x_1,\ldots,x_m\in\RR^n\) их выпуклая комбинация имеет вид
\[
    z=\sum_{i=1}^{m}\theta_i x_i,
    \qquad
    \theta_i\ge0,
    \qquad
    \sum_{i=1}^{m}\theta_i=1.
\]



\begin{definition}[Выпуклое множество]
Множество \(C\subseteq\RR^n\) называется выпуклым, если
\[
    x,y\in C,\quad \theta\in[0,1]
    \quad\Longrightarrow\quad
    \theta x+(1-\theta)y\in C.
\]
\end{definition}

Иными словами, вместе с любыми двумя точками выпуклое множество
содержит весь соединяющий их отрезок.

Интервалы, аффинные подпространства и евклидовы шары являются
выпуклыми. Отсутствие разрывов или отверстий недостаточно: множество
выпукло именно тогда, когда содержит отрезок между каждой парой своих
точек. Например,
\[
    C=\{2,5\}
\]
не выпукло, поскольку
\[
    \frac12\cdot2+\frac12\cdot5=3.5\notin C.
\]

\begin{definition}[Выпуклая оболочка]
Выпуклой оболочкой множества \(S\subseteq\RR^n\) называется множество
всех конечных выпуклых комбинаций его точек:
\[
    \operatorname{conv}(S)
    =
    \left\{
        \sum_{i=1}^{m}\theta_i x_i
        \,\middle|\,
        m\in\mathbb N,\quad
        x_i\in S,\quad
        \theta_i\ge0,\quad
        \sum_{i=1}^{m}\theta_i=1
    \right\}.
\]
\end{definition}

\section{Выпуклые функции и минимумы}

Пусть \(C\subseteq\RR^n\) --- выпуклое множество.

\begin{definition}[Выпуклая функция]
Функция \(f\colon C\to\RR\) называется выпуклой, если для любых
\(x,y\in C\) и \(\theta\in[0,1]\)
\[
    f\bigl(\theta x+(1-\theta)y\bigr)
    \le
    \theta f(x)+(1-\theta)f(y).
\]
\end{definition}

Правая часть задаёт значение хорды между точками
\[
    (x,f(x))
    \quad\text{и}\quad
    (y,f(y)).
\]
Следовательно, граф выпуклой функции лежит не выше любой своей хорды.

Требование выпуклости области определения существенно: точка
\[
    \theta x+(1-\theta)y
\]
должна принадлежать области определения функции.

\begin{definition}[Строго выпуклая функция]
Функция \(f\colon C\to\RR\) называется строго выпуклой, если для
различных \(x,y\in C\) и каждого \(\theta\in(0,1)\)
\[
    f\bigl(\theta x+(1-\theta)y\bigr)
    <
    \theta f(x)+(1-\theta)f(y).
\]
\end{definition}

Постоянная функция выпукла, но не строго выпукла. Выпуклая функция
может иметь несколько минимумов, тогда как строго выпуклая функция
имеет не более одного глобального минимума.



\begin{theorem}
Пусть \(C\) выпукло, а \(f\colon C\to\RR\) выпукла. Тогда всякий
локальный минимум \(f\) является глобальным.
\end{theorem}

\begin{proof}
Пусть \(x^\star\) --- локальный минимум относительно множества
\(C\), но существует \(y\in C\),
для которого
\[
    f(y)<f(x^\star).
\]
Для \(\theta\in(0,1)\) положим
\[
    x_\theta=(1-\theta)x^\star+\theta y.
\]
При достаточно малом \(\theta\) точка \(x_\theta\) сколь угодно
близка к \(x^\star\). По выпуклости
\[
    f(x_\theta)
    \le
    (1-\theta)f(x^\star)+\theta f(y)
    <
    f(x^\star),
\]
что противоречит локальной минимальности \(x^\star\).
\end{proof}

\begin{corollary}
Если строго выпуклая функция достигает минимума, то этот минимум
единственен.
\end{corollary}

Для общей невыпуклой функции локальные минимумы могут иметь различные
значения. Поэтому локальная сходимость алгоритма не обеспечивает
глобальной оптимальности.



Функции
\[
    f(x)=\frac12x^2,
    \qquad
    f(x)=\frac14x^4,
    \qquad
    f(x)=\mathrm{e}^x
\]
выпуклы на \(\RR\).

Функция
\[
    f(x)=c
\]
одновременно выпукла и вогнута.

Функция
\[
    f(x)=\log x,
    \qquad x>0,
\]
не выпукла, но вогнута. Поэтому
\[
    -\log x
\]
выпукла на \(\RR_{++}\).

Функция может быть одновременно невыпуклой и невогнутой. Поэтому
невыпуклость не означает вогнутость.

\begin{definition}[Вогнутая функция]
Функция \(f\) называется вогнутой, если функция \(-f\) выпукла.
Эквивалентно,
\[
    f\bigl(\theta x+(1-\theta)y\bigr)
    \ge
    \theta f(x)+(1-\theta)f(y).
\]
\end{definition}

После смены знака минимизация вогнутой функции превращается в
максимизацию выпуклой функции, которая в общем случае не является
задачей выпуклой оптимизации. Стандартные выпуклые постановки ---
минимизация выпуклой или максимизация вогнутой функции на выпуклом
множестве.

\section{Операции, сохраняющие выпуклость}



\begin{theorem}
Пусть \(f_1,\ldots,f_m\) выпуклы на общем выпуклом множестве \(C\), а
\[
    \alpha_i\ge0.
\]
Тогда функция
\[
    f(x)=\sum_{i=1}^{m}\alpha_i f_i(x)
\]
выпукла на \(C\).
\end{theorem}

\begin{proof}
Для \(x,y\in C\) и \(\theta\in[0,1]\)
\[
\begin{aligned}
    f\bigl(\theta x+(1-\theta)y\bigr)
    &=
    \sum_{i=1}^{m}
    \alpha_i
    f_i\bigl(\theta x+(1-\theta)y\bigr)
    \\
    &\le
    \sum_{i=1}^{m}
    \alpha_i
    \bigl(
        \theta f_i(x)+(1-\theta)f_i(y)
    \bigr)
    \\
    &=
    \theta f(x)+(1-\theta)f(y).
\end{aligned}
\]
\end{proof}

Поскольку функция \(t\mapsto\lvert t\rvert\) выпукла,
\[
    \lVert x\rVert_1
    =
    \sum_{i=1}^{n}\lvert x_i\rvert
\]
также выпукла.

Аналогично функция
\[
    f(x)
    =
    \frac12x^2+\frac14x^4
\]
выпукла как сумма выпуклых функций.



\begin{theorem}
Пусть \(f_1,\ldots,f_m\) выпуклы на \(C\). Тогда функция
\[
    f(x)=\max_{1\le i\le m}f_i(x)
\]
выпукла на \(C\).
\end{theorem}

\begin{proof}
Для каждого \(i\)
\[
\begin{aligned}
    f_i\bigl(\theta x+(1-\theta)y\bigr)
    &\le
    \theta f_i(x)+(1-\theta)f_i(y)
    \\
    &\le
    \theta f(x)+(1-\theta)f(y).
\end{aligned}
\]
Переход к максимуму по \(i\) даёт требуемое неравенство.
\end{proof}

Поточечные максимумы возникают в негладких целевых функциях,
минимаксных постановках и задачах робастной оптимизации.



\begin{theorem}
Пусть \(D\subseteq\RR^m\) выпукло, \(f\colon D\to\RR\) выпукла,
\(A\in\RR^{m\times n}\) и \(b\in\RR^m\). Тогда множество
\[
    C=\{x\in\RR^n:Ax+b\in D\}
\]
выпукло, а функция
\[
    g(x)=f(Ax+b)
\]
выпукла на \(C\).
\end{theorem}

\begin{proof}
Для \(x,y\in C\) и \(\theta\in[0,1]\)
\[
    A\bigl(\theta x+(1-\theta)y\bigr)+b
    =
    \theta(Ax+b)+(1-\theta)(Ay+b)\in D,
\]
поэтому \(C\) выпукло. Кроме того,
\[
\begin{aligned}
    g\bigl(\theta x+(1-\theta)y\bigr)
    &=
    f\bigl(
        \theta(Ax+b)+(1-\theta)(Ay+b)
    \bigr)
    \\
    &\le
    \theta f(Ax+b)+(1-\theta)f(Ay+b).
\end{aligned}
\]
\end{proof}

В частности, поскольку квадрат евклидовой нормы выпуклый,
\[
    f(x)=\frac12\lVert Ax-b\rVert_2^2
\]
является выпуклой функцией.

\section{Логарифмический барьер}

Для строгого линейного ограничения
\[
    a^{\mathsf T}x<b
\]
можно использовать барьер
\[
    \phi(x)
    =
    -\log\bigl(b-a^{\mathsf T}x\bigr).
\]
Функция определена только внутри допустимой области. При приближении
к её границе
\[
    b-a^{\mathsf T}x\to0^+
\]
выполняется
\[
    \phi(x)\to+\infty.
\]

Для ограничений \(a_i^{\mathsf T}x\le b_i\) барьерный метод
рассматривает семейство задач
\[
    \min_{x:\,a_i^{\mathsf T}x<b_i}
    \left\{
        f(x)
        -
        \mu
        \sum_{i=1}^{m}
        \log\bigl(b_i-a_i^{\mathsf T}x\bigr)
    \right\},
    \qquad
    \mu>0.
\]
При уменьшении \(\mu\) решения барьерных задач приближаются к решению
исходной задачи при стандартных условиях регулярности. Поскольку
\(-\log t\) выпукла на \(t>0\), а аффинная композиция сохраняет
выпуклость, каждый барьер выпуклый на своей области определения.

\section{Условия выпуклости первого и второго порядка}

\begin{theorem}[Критерий первого порядка]
Пусть \(C\subseteq\RR^n\) открыто и выпукло, а
\(f\colon C\to\RR\) непрерывно дифференцируема. Тогда \(f\) выпукла
на \(C\) тогда и только тогда, когда
\[
    f(y)
    \ge
    f(x)
    +
    \langle\nabla f(x),y-x\rangle
\]
для всех \(x,y\in C\).
\end{theorem}

Геометрически это означает, что граф выпуклой дифференцируемой
функции лежит не ниже любой касательной гиперплоскости.

\begin{proof}
Пусть сначала \(f\) выпукла. Для \(\lambda\in(0,1]\)
\[
\begin{aligned}
    f\bigl(x+\lambda(y-x)\bigr)
    &=
    f\bigl((1-\lambda)x+\lambda y\bigr)
    \\
    &\le
    (1-\lambda)f(x)+\lambda f(y).
\end{aligned}
\]
Следовательно,
\[
    \frac{
        f\bigl(x+\lambda(y-x)\bigr)-f(x)
    }{\lambda}
    \le
    f(y)-f(x).
\]
Переход к пределу при \(\lambda\to0^+\) даёт
\[
    \langle\nabla f(x),y-x\rangle
    \le
    f(y)-f(x).
\]

Обратно, пусть неравенство первого порядка выполнено. Для
\(z=\theta x+(1-\theta)y\) применим его к парам \((z,x)\) и
\((z,y)\):
\[
\begin{aligned}
    f(x)&\ge f(z)+\langle\nabla f(z),x-z\rangle,\\
    f(y)&\ge f(z)+\langle\nabla f(z),y-z\rangle.
\end{aligned}
\]
Умножая неравенства на \(\theta\) и \(1-\theta\) и складывая,
получаем
\[
    \theta f(x)+(1-\theta)f(y)\ge f(z),
\]
поскольку
\(\theta(x-z)+(1-\theta)(y-z)=0\). Следовательно, \(f\) выпукла.
\end{proof}

\begin{corollary}[Условие глобального минимума]
Пусть \(C\) открыто и выпукло, \(f\colon C\to\RR\) выпукла и
дифференцируема. Если \(x^\star\in C\) и
\[
    \nabla f(x^\star)=0,
\]
то \(x^\star\) является глобальным минимумом \(f\) на \(C\).
\end{corollary}

\begin{proof}
По критерию первого порядка для любого \(y\in C\)
\[
    f(y)
    \ge
    f(x^\star)
    +
    \langle\nabla f(x^\star),y-x^\star\rangle
    =
    f(x^\star).
\]
\end{proof}



\begin{theorem}[Критерий второго порядка]
Пусть \(C\subseteq\RR^n\) открыто и выпукло, а
\(f\colon C\to\RR\) дважды непрерывно дифференцируема. Тогда \(f\)
выпукла на \(C\) тогда и только тогда, когда
\[
    \nabla^2f(x)\succeq0
    \qquad
    \text{для всех }x\in C.
\]
\end{theorem}

\begin{proof}
Если \(f\) выпукла, то для любых \(x\in C\) и \(v\in\RR^n\)
одномерная функция \(g(t)=f(x+tv)\) выпукла в окрестности нуля.
Следовательно,
\[
    v^{\mathsf T}\nabla^2f(x)v=g''(0)\ge0.
\]

Обратно, пусть \(\nabla^2f(x)\succeq0\) на \(C\). Для любых
\(x,y\in C\) функция
\[
    g(t)=f\bigl(x+t(y-x)\bigr),
    \qquad t\in[0,1],
\]
удовлетворяет
\[
    g''(t)
    =
    (y-x)^{\mathsf T}
    \nabla^2f\bigl(x+t(y-x)\bigr)
    (y-x)
    \ge0.
\]
Поэтому \(g\) выпукла на \([0,1]\), что равносильно выпуклости
\(f\) на каждом отрезке в \(C\).
\end{proof}

Положительная полуопределённость означает
\[
    v^{\mathsf T}\nabla^2f(x)v\ge0
    \qquad
    \text{для всех }v\in\RR^n.
\]

В одномерном случае критерий принимает вид
\[
    f''(x)\ge0.
\]

Для функции
\[
    f(x)=\frac12x^2
\]
имеем
\[
    f''(x)=1>0,
\]
поэтому она выпукла.

Для функции
\[
    f(x)=\frac14x^4
\]
получаем
\[
    f''(x)=3x^2\ge0.
\]
Следовательно, она также выпукла, хотя
\[
    f''(0)=0.
\]

\section{Схема проверки выпуклости}

\begin{algorithm}[H]
\caption{Проверка выпуклости функции}
\label{alg:convexity-check}
\begin{algorithmic}[1]

\Require
функция \(f\colon C\to\RR\)

\State
проверить выпуклость области \(C\)

\If{\(f\) дважды непрерывно дифференцируема}

    \State
    вычислить \(\nabla^2f(x)\)

    \State
    проверить
    \(\nabla^2f(x)\succeq0\) для всех \(x\in C\)

\ElsIf{\(f\) непрерывно дифференцируема}

    \State
    проверить
    \[
        f(y)
        \ge
        f(x)+
        \langle\nabla f(x),y-x\rangle
    \]
    для всех \(x,y\in C\)

\Else

    \State
    проверить определение выпуклости непосредственно

\EndIf

\Ensure
заключение о выпуклости \(f\)

\end{algorithmic}
\end{algorithm}


Множество выпукло, если оно содержит отрезок между любыми двумя
своими точками. На выпуклой области всякий локальный минимум выпуклой
функции является глобальным. Неотрицательные суммы, поточечные
максимумы и аффинные композиции сохраняют выпуклость. Для непрерывно
дифференцируемой функции выпуклость эквивалентна условию
\[
    f(y)
    \ge
    f(x)+\langle\nabla f(x),y-x\rangle,
\]
а для дважды непрерывно дифференцируемой функции --- условию
\[
    \nabla^2f(x)\succeq0.
\]